% Part: turing-machines
% Chapter: undecidability
% Section: trakhtenbrot

\documentclass[../../../include/open-logic-section]{subfiles}

\begin{document}

\olfileid{tur}{und}{tra}
\olsection{ట్రాఖ్టెన్‌బ్రోట్ సిద్ధాంతం}

\begin{explain}
  \olref[rep]{sec}లో ట్యూరింగ్ యంత్రం~$M$, ఇన్‌పుట్ సంకేతమాల~$w$లకు
  \tetoken{వాక్యాలు}{!!{sentence}s} $!T(M,w)$,~$!E(M,w)$లను నిర్వచించాం. తరువాత
  \olref[ver]{lem:valid-if-halt}, \olref[ver]{lem:halt-if-valid}లలో,
  ఇన్‌పుట్~$w$పై ప్రారంభించిన $M$ చివరకు ఆగితే, అప్పుడే మాత్రమే
  $!T(M,w) \lif !E(M,w)$ చెల్లుబాటవుతుందని చూపాం. ఆగే సమస్య
  అనిర్ణయనీయమైనది కనుక, మొదటిస్థాయి తర్కపు \tetoken{వాక్యాల}{!!{sentence}s} చెల్లుబాటుతనం,
  సంతృప్తిపరచదగినత అనిర్ణయనీయాలని ఇది సూచిస్తుంది
  (\cref{tur:und:uns:thm:decision-prob,tur:und:uns:cor:undecidable-sat}).

  కానీ వాక్యాల చెల్లుబాటుతనం, సంతృప్తిపరచదగినత ఏకపక్ష \tetoken{నిర్మాణాలకు}{!!{structure}s},
  అవి పరిమితమైనా అనంతమైనా, నిర్వచించబడతాయి. \tetoken{ఒక వాక్యం}{!!a{sentence}} పరిమిత
  \tetoken{నిర్మాణంలో}{!!{structure}} సంతృప్తిపరచదగినదో (లేదా అన్ని పరిమిత \tetoken{నిర్మాణాలలో}{!!{structure}s}
  చెల్లుబాటవుతుందో) నిర్ణయించడం సులభమని మీరు అనుమానించవచ్చు. అలా కాదని
  చూపేలా నిర్ణయ సమస్య అపరిష్కార్యత నిరూపణను అనుకూలించవచ్చు.

  మొదట, \olref[ver]{lem:halt-if-valid} నిరూపణకు తిరిగి వెళితే, ఇన్‌పుట్~$w$పై
  ప్రారంభించిన యంత్రం~$M$ కచ్చితంగా చేసే పనిని వర్ణించే $!T(M,w)$ నమూనా
  $\Struct M$ను అక్కడ నిర్మించామని చూస్తారు. ఆ నమూనా వ్యక్తి క్షేత్రం
  $\Nat$, అంటే అనంతం. కానీ $M$ ఇన్‌పుట్~$w$పై నిజంగా ఆగితే, అదే విధంగా
  పరిమిత నమూనా~$\Struct M'$ను నిర్మించవచ్చు. ఇన్‌పుట్~$w$పై ప్రారంభించిన
  $M$, $k$ అంచెల తరువాత ఆగుతుందని అనుకుందాం. వ్యక్తి క్షేత్రం కోసం $n$ను
  $k$ కంటే ఒకటి ఎక్కువ, $w$ పొడవు---వాటిలో పెద్దదిగా ఉంచి,
  $\{0,\dots,n\}$ సమితిని వ్యక్తి క్షేత్రం~$\Domain{M'}$గా తీసుకోండి;
  ఇంకా
  \[
    \Assign{\prime}{M'}(x) =
    \begin{cases}
      x + 1 &\text{$x < n$ అయితే}\\
      n &\text{లేకపోతే,}
    \end{cases}
  \]
  అని, $\tuple{x,y}\in\Assign{<}{M'}$ అయితే, అప్పుడే మాత్రమే $x<y$ లేదా
  $x=y=n$ అని ఉంచండి. మిగతా విధంగా $\Struct{M'}$ను $\Struct{M}$లాగే
  నిర్వచిస్తాం. $\Struct{M'}$ నిర్వచనం వల్ల,
  \olref[ver]{lem:halt-if-valid} నిరూపణలో మాదిరిగానే,
  $\Sat{M'}{!T(M,w)}$. అలాగే $M$ ఇన్‌పుట్~$w$పై ఆగుతుందని అనుకున్నాం కనుక,
  $\Sat{M'}{!E(M,w)}$. అందువల్ల $\Struct{M'}$ అనేది
  $!T(M,w)\land !E(M,w)$కు ఒక పరిమిత నమూనా (ఇక్కడ $\lif$ను~$\land$తో
  మార్చామని గమనించండి).\sourcecorrection{OLTETURUND-014}{పరిమిత నమూనా వ్యక్తి క్షేత్రానికి ఆంగ్ల మూలం $n$ను $k$, $\len w$లలో పెద్దదిగా మాత్రమే తీసుకుంది. ప్రారంభ శీర్ష స్థానం~$1$ కావడం వల్ల $k$ అంచెలలో శీర్షం $k+1$ వరకు చేరవచ్చు; అంతేకాక తరువాత చేర్చే $!B$ షరతు చివరి నిజమైన అంచెను పై మూలకంలోకి మడవకూడదు. అందువల్ల అవసరమైన సరిహద్దు $n=\max(k+1,\len w)$గా ఇచ్చాం; ఇదే మరమ్మతును కింది నిర్మాణంలోనూ వాడాం.}

  నిరూపణలో సగం పూర్తయింది: $M$ ఇన్‌పుట్~$w$పై ఆగితే,
  $!T(M,w)\land !E(M,w)$కు పరిమిత నమూనా ఉంటుందని చూపాం. దురదృష్టవశాత్తూ
  దీని విపర్యయం నిజం కాదు; అంటే, కొన్ని ఇన్‌పుట్~$w$పై ఆగని ట్యూరింగ్
  యంత్రాలు ఉన్నా, $!T(M,w)\land !E(M,w)$కు పరిమిత నమూనా ఉండవచ్చు.
  ఉదాహరణకు, ఒకే స్థితి $q_0$, ఒకే నిర్దేశం
  $\delta(q_0,\TMblank)=\tuple{q_0,\TMblank,\TMstay}$ గల యంత్రం~$M$ను
  పరిశీలించండి. ఖాళీ ఇన్‌పుట్~$w=\emptyseq$పై ప్రారంభించినప్పుడు ఈ యంత్రం
  ఎప్పటికీ ఆగదు: అది అనంత చక్రంలో ఉంటుంది, కానీ టేపును మార్చదు, శీర్షాన్ని
  కదపదు. అన్ని స్థితివిన్యాసాలూ ఒకటే (అదే స్థితి, అదే శీర్ష స్థానం, అదే
  టేపు విషయం). $!T(M,\emptyseq)\land !E(M,\emptyseq)$ను సంతృప్తిపరిచే
  పరిమిత \tetoken{నిర్మాణం}{!!{structure}}~$\Struct{M''}$ను నిర్వచించవచ్చు (అభ్యాసం). అయితే,
  అలాంటి \tetoken{నిర్మాణాలు}{!!{structure}s} తొలగిపోయేలా $!T(M,w)$ను తగిన విధంగా మార్చవచ్చు.

  \begin{prob}
    $M$ ఒకే స్థితి $q_0$, ఒకే నిర్దేశం
    $\delta(q_0,\TMblank)=\tuple{q_0,\TMblank,\TMstay}$ గల ట్యూరింగ్ యంత్రం
    అనుకుందాం. $\Domain{M''}=\{0,1,2\}$,
    $\Assign{\prime}{M''}(0)=\Assign{\prime}{M''}(1)=1$ మరియు
    $\Assign{\prime}{M''}(2)=2$ అని, ఇంకా
    $\Assign{<}{M''}=\{\tuple{0,1},\tuple{1,1},\tuple{2,2}\}$ అని ఉంచండి.
    $!T(M,\emptyseq)$, $!E(M,\emptyseq)$ సత్యమయ్యేలా
    $\Assign{\Obj Q_{q_0}}{M''}$, $\Assign{\Obj S_{\TMblank}}{M''}$,
    $\Assign{\Obj S_{\TMendtape}}{M''}$లను నిర్వచించి, అవి ఎందుకు సత్యమో
    వివరించండి. సూచన: $\delta(q_0,\TMendtape)$ నిర్వచితం కాదని గమనించండి.
    కింది రెండూ $\Struct{M''}$లో సత్యమయ్యేలా చూడండి:
    \begin{align*}
    & \Obj Q_{q_0}(\num{1}, \num{n}) \land \Obj S_{\TMendtape}(\num{0}, \num{n})
      \land
      \lforall[x][(\num{0} < x \lif \Obj S_\TMblank(x, \num{n}))] \text{\quad ప్రతి $n \in \Nat$కు}\\
    & \lexists[y][(\Obj Q_{q_0}(\num{0}, y) \land \Obj S_{\TMendtape}(\num{0}, y))]
    \end{align*}
  \end{prob}\sourcecorrection{OLTETURUND-015}{ఈ అభ్యాసంలోని ఆంగ్ల మూలం ఏకైక నిర్దేశపు ఫలిత స్థితిని నిర్వచించని $q$గా వ్రాసి, ఉత్తరవర్తి అర్థనిర్దేశంలో $\Assign{\prime}{M''}(1)$కు బదులు వేరే నిర్మాణం~$M'$ను పేర్కొంది. ఉదాహరణ ముందు వాక్యానికీ ఒకే స్థితి అన్న షరతుకూ అనుగుణంగా $q_0$ను, అభ్యాసపు నిర్మాణం అంతటా $M''$ను పునరుద్ధరించాం.}
\end{explain}

ఇన్‌పుట్ $w=\sigma_{i_1}\dots\sigma_{i_k}$పై ట్యూరింగ్ యంత్రం~$M$ చర్యను
వర్ణించే \tetoken{వాక్యాలను}{!!{sentence}s} పరిశీలించండి:
\begin{enumerate}
\item సంఖ్యలనూ~$<$నూ వర్ణించే స్వీకృతాలు
  (\olref[rep]{sec}లోని $!T(M,w)$ నిర్వచనంలో మాదిరిగానే).
\item ఇన్‌పుట్ స్థితివిన్యాసాన్ని వర్ణించే స్వీకృతాలు: $!T(M,w)$ నిర్వచనంలో
  ఉన్నట్లే.
\item ఒక స్థితివిన్యాసం నుండి తరువాతిదానికి సంక్రమణను వర్ణించే స్వీకృతాలు:

కింది వాటి కోసం, $!A(x,y)$ను మునుపటిలాగే ఉంచి, ఇంకా
\[
  !B(y) \ident \lforall[x][(x < y \lif \eq/[x][y])].
\]
అని ఉంచుదాం.
\begin{enumerate}
\item \ollabel{rep-right} ప్రతి నిర్దేశం $\delta(q_i,\sigma)=
  \tuple{q_j,\sigma',\TMright}$కు కింది \tetoken{వాక్యం}{!!{sentence}}:
\begin{align*}
& \lforall[x][\lforall[y][(
   (\Obj Q_{q_i}(x, y) \land \Obj S_{\sigma}(x, y)) \lif {}]] \\
&\qquad   (\Obj Q_{q_j}(x', y') \land \Obj S_{\sigma'}(x, y') \land
!A(x, y) \land !B(y')))
\end{align*}
\item \ollabel{rep-left} ప్రతి నిర్దేశం $\delta(q_i,\sigma)=
  \tuple{q_j,\sigma',\TMleft}$కు కింది \tetoken{వాక్యం}{!!{sentence}}:
\begin{align*}
& \lforall[x][\lforall[y][
    ((\Obj Q_{q_i}(x', y) \land \Obj S_{\sigma}(x', y)) \lif {}]]\\
& \qquad   (\Obj Q_{q_j}(x, y') \land \Obj S_{\sigma'}(x', y') \land
!A(x', y) \land !B(y'))) \land {}\\
& \lforall[y][((\Obj Q_{q_i}(\num{0}, y) \land \Obj S_{\sigma}(\num{0},
    y)) \lif {}]\\
& \qquad (\Obj Q_{q_j}(\num{0}, y') \land \Obj S_{\sigma'}(\num{0},
  y') \land !A(\num{0}, y) \land !B(y')))
\end{align*}
\item \ollabel{rep-stay} ప్రతి నిర్దేశం $\delta(q_i,\sigma)=
  \tuple{q_j,\sigma',\TMstay}$కు కింది \tetoken{వాక్యం}{!!{sentence}}:
\begin{align*}
& \lforall[x][\lforall[y][(
   (\Obj Q_{q_i}(x, y) \land \Obj S_{\sigma}(x, y)) \lif {}]] \\
&\qquad (\Obj Q_{q_j}(x, y') \land \Obj S_{\sigma'}(x, y') \land
!A(x, y) \land !B(y')))
\end{align*}
\end{enumerate}
చూడగలిగినట్లుగా, $M$ సంక్రమణలను వర్ణించే \tetoken{వాక్యాలు}{!!{sentence}s} $!T(M,w)$లోని
అనుగుణ \tetoken{వాక్యంతో}{!!{sentence}} ఒకటే; చివర్లో $!B(y')$ను చేర్చడం మాత్రమే తేడా.
$!B(y')$ అనేది ``తరువాతి'' స్థితివిన్యాసపు సంఖ్య $y'$ ముందరి సంఖ్యలైన
$\Obj 0$, $\Obj 0'$, \dots{} అన్నింటి నుండీ వేరుగా ఉండేలా చూస్తుంది.\sourcecorrection{OLTETURUND-016}{మార్చిన ఎడమ-కదలిక స్వీకృతంలో ఆంగ్ల మూలం సున్నా కాని గడి శాఖకు $!B(y')$ను చేర్చలేదు; అందువల్ల ఎడమకు కదులుతూ ఎప్పటికీ ఆగని నడకకు కొత్త సమయ విలువలను బలవంతం చేయదు, కింది అనంతత్వ వాదన విఫలమవుతుంది. అదే శాఖ $!A(x,y)$తో వ్రాసిన గడికి బదులు గమ్య గడిని మినహాయించే OLTETURUND-006 దోషాన్నీ పునరావృతం చేసింది. అన్ని మూడు కదలిక రకాలకూ $!B(y')$ను ఏకరీతిగా విధించి, ఎడమ శాఖలో $!A(x',y)$ను వాడాం.}
% \item సంగతత్వ షరతులను వ్యక్తం చేసే వాక్యాలు:
% \begin{enumerate}
%   \item\ollabel{rep-con-symb}%
%   ఏ గడిపైనా ఒకే సమయంలో ఒకటికంటే ఎక్కువ సంకేతాలు ఉండలేవని చెప్పే
%   \tetoken{ఒక వాక్యం}{!!^a{sentence}}:
%   \[\lforall[x][\lforall[y][(\Obj S_{\sigma}(x, y) \lif \lnot \Obj
%   S_{\sigma'}(x, y))]],\]
%   $T$ యొక్క ప్రతి రెండు వేర్వేరు టేపు సంకేతాలు $\sigma\neq\sigma'$కు.
%   \item\ollabel{rep-con-state}%
%   ఏ సమయంలోనూ $M$ ఒకటికంటే ఎక్కువ స్థితుల్లో ఉండలేదని చెప్పే
%   \tetoken{ఒక వాక్యం}{!!^a{sentence}}:
%   \[\lforall[x][\lforall[y][(\Obj Q_{q}(x, y) \lif
%   \lforall[z][\lnot \Obj
%   Q_{q'}(z, y))]],\]
%   $T$ యొక్క ప్రతి రెండు వేర్వేరు స్థితులు $q\neq q'$కు.
%   \item\ollabel{rep-con-tape}%
%   ఏ సమయంలోనూ $M$ ఒకటికంటే ఎక్కువ టేపు గడులపై ఉండలేదని చెప్పే
%   \tetoken{ఒక వాక్యం}{!!^a{sentence}}:
%   \[\lforall[x][\lforall[y][(\Obj Q_{q}(x, y) \lif
%   \lforall[z][\eq/[z][y]\lif \lnot \Obj
%   Q_{q'}(z, y))]],\]
%   $T$ యొక్క ప్రతి స్థితుల యుగ్మం $q$, $q'$కు.
% \end{enumerate}
\end{enumerate}
పై \tetoken{వాక్యాలన్నిటి}{!!{sentence}s} సంయోగాన్ని ట్యూరింగ్ యంత్రం~$M$, ఇన్‌పుట్~$w$లకు
$!T'(M,w)$గా ఉంచుదాం.

\begin{lem}\ollabel{lem:halts-sat}
  ఇన్‌పుట్~$w$పై ప్రారంభించిన $M$ ఆగితే,
  $!T'(M,w)\land !E(M,w)$కు పరిమిత నమూనా ఉంటుంది.
\end{lem}

\begin{proof}
  $\Struct{M'}$ను \olref[ver]{lem:halt-if-valid} నిరూపణలో ఉన్నట్లే
  తీసుకుందాం; అయితే
  \begin{align*}
    \Domain{M'} & = \{0, \dots, n\},\\
    \Assign{\prime}{M'}(x) & =
    \begin{cases}
      x + 1 &\text{$x < n$ అయితే}\\
      n &\text{లేకపోతే,}
    \end{cases} \\
    \tuple{x,y} \in \Assign{<}{M'} &\text{ అయితే, అప్పుడే మాత్రమే $x < y$ లేదా $x = y = n$,}
  \end{align*}
  ఇక్కడ $n=\max(k+1,\len{w})$; $k$ అనేది ఇన్‌పుట్~$w$పై ప్రారంభించిన
  $M$ $k$ అంచెల తరువాత ఆగే అత్యల్ప సంఖ్య.
  $\Sat{M'}{!T'(M,w)\land !E(M,w)}$ అని నిర్ధారించడాన్ని అభ్యాసంగా
  వదిలేస్తాం.
\end{proof}

\begin{prob}
  $\Sat{M'}{!T'(M,w)\land !E(M,w)}$ అని నిరూపించి
  \olref[tur][und][tra]{lem:halts-sat} నిరూపణను పూర్తి చేయండి.
\end{prob}\sourcecorrection{OLTETURUND-017}{ఈ ఉపసిద్ధాంత నిరూపణలో ఆంగ్ల మూలం $!E$లోని ఆశ్చర్యార్థక సంకేతాన్ని వదిలింది; వెంటనున్న అభ్యాసం మరింతగా సవరణకు ముందు ఉన్న $!T$నే అడిగి, $!T'$ను వదిలి, $!E$ ఆశ్చర్యార్థక సంకేతాన్నీ వదిలింది. ఉపసిద్ధాంతం నిజంగా పేర్కొన్న సూత్రమైన $!T'(M,w)\land !E(M,w)$ను రెండుచోట్లా పునరుద్ధరించాం; OLTETURUND-014 సరిహద్దు $n=\max(k+1,\len w)$ను కూడా ఇక్కడ ఏకరీతిగా వాడాం.}

\begin{lem}\ollabel{lem:sat-halts}
  $!T'(M,w)\land !E(M,w)$కు పరిమిత నమూనా ఉంటే, ఇన్‌పుట్~$w$పై
  ప్రారంభించిన $M$ ఆగుతుంది.
\end{lem}

\begin{proof}
  వ్యతిరేకాన్ని చూపుతాం. ఇన్‌పుట్~$w$పై ప్రారంభించిన $M$ ఆగదని అనుకుందాం.
  $!T'(M,w)\land !E(M,w)$కు అసలు నమూనానే లేకపోతే, మన పని పూర్తయింది.
  కాబట్టి $\Struct{M}$ అనేది $!T'(M,w)\land !E(M,w)$కు ఒక నమూనా
  అనుకుందాం. అది పరిమితం కాలేదని చూపాలి.

  \olref[ver]{lem:config}లో మాదిరిగానే, ఇన్‌పుట్~$w$పై ప్రారంభించిన $M$
  $n$వ అంచెకు ముందే ఆగకపోతే
  $!T'(M,w)\Entails !C(M,w,n)$ అని, ఇంకా $n>0$ అయితే
  $!T'(M,w)\Entails !B(\num{n})$ అని నిరూపించవచ్చు. ఇన్‌పుట్~$w$పై
  ప్రారంభించిన $M$ ఆగదు కనుక, ఇవి తగిన అన్ని $n\in\Nat$లకు వర్తిస్తాయి.
  \olref[rep]{prop:mlessk}లోని అదే సంఖ్య-స్వీకృత వాదన వల్ల, ప్రతి $k<n$కు
  $!T'(M,w)\Entails \num{k}<\num{n}$ అని గమనించండి. అలాగే
  $!B(\num{n})\Entails
  \num{k}<\num{n}\lif\eq/[\num{k}][\num{n}]$. కాబట్టి ప్రతి $k<n$కు
  $\Sat{M}{\eq/[\num{k}][\num{n}]}$; అంటే అనంతంగా ఉన్న
  $\num{k}$ పదాలన్నీ $\Struct{M}$లో వేర్వేరు విలువలు కలిగి ఉండాలి. కానీ
  దీనికి $\Domain{M}$ అనంతంగా ఉండాలి. అందువల్ల $\Struct{M}$ అనేది
  $!T'(M,w)\land !E(M,w)$కు పరిమిత నమూనా కాలేదు.
\end{proof}\sourcecorrection{OLTETURUND-018}{ఈ వ్యతిరేక నిరూపణలో ఆంగ్ల మూలం ఉపసిద్ధాంత పరికల్పనలోని $!T'$కు బదులు సవరణకు ముందు ఉన్న $!T$ నమూనాను అనుకుంది; అలా అయితే $!B$ షరతులు లభించవు. నమూనా పరికల్పనను $!T'(M,w)\land !E(M,w)$గా పునరుద్ధరించాం. అదనంగా, మూలం $!B(\num n)$ను $n=0$ సహా అన్నిటికీ పేర్కొంది; కానీ $!B(\num0)$ ఏ సంక్రమణ ఫలితభాగం నుండీ రాదు, అనంతత్వ వాదనకు అవసరమూ కాదు. $!B(\num n)$ దావాను $n>0$కు కచ్చితంగా పరిమితం చేశాం.}

\begin{prob}
  ఇన్‌పుట్~$w$పై ప్రారంభించిన $M$ $n$వ అంచెకు ముందే ఆగకపోతే, ప్రతి
  $n>0$కు $!T'(M,w)\Entails !B(\num{n})$ అని నిరూపించి
  \olref[tur][und][tra]{lem:sat-halts} నిరూపణను పూర్తి చేయండి.
\end{prob}

\begin{thm}[ట్రాఖ్టెన్‌బ్రోట్ సిద్ధాంతం]
  \ollabel{thm:trakhtenbrodt}
  మొదటిస్థాయి తర్కంలోని ఏకపక్ష \tetoken{వాక్యానికి}{!!{sentence}} పరిమిత నమూనా ఉందో లేదో
  (అంటే, అది పరిమితంగా సంతృప్తిపరచదగినదో కాదో) నిర్ణయించలేం.
\end{thm}

\begin{proof}
  పరిమిత సంతృప్తిపరచదగినత సమస్యను నిర్ణయించే ట్యూరింగ్ యంత్రం~$F$ ఉందని
  అనుకుందాం. అప్పుడు ఏ ట్యూరింగ్ యంత్రం~$M$, ఇన్‌పుట్~$w$ ఇచ్చినా,
  $!T'(M,w)\land !E(M,w)$ అనే వాక్యాన్ని గణించి, దానికి పరిమిత నమూనా
  ఉందో లేదో నిర్ణయించడానికి~$F$ను ఉపయోగించగలిగేవారం.
  \cref{tur:und:tra:lem:halts-sat,tur:und:tra:lem:sat-halts} ప్రకారం,
  ఇన్‌పుట్~$w$పై ప్రారంభించిన $M$ ఆగితే, అప్పుడే మాత్రమే దానికి పరిమిత
  నమూనా ఉంటుంది. కాబట్టి ఆగే సమస్యను పరిష్కరించడానికి $F$ను ఉపయోగించగలిగేవారం;
  అది పరిష్కరించలేనిదని మనకు తెలుసు.
\end{proof}

\begin{cor}\ollabel{cor:fproof-incomp}%
  \emph{పరిమిత} చెల్లుబాటుతనానికి నిర్దుష్టమైనదీ సంపూర్ణమైనదీ అయిన
  \tetoken{వ్యుత్పత్తి}{!!{derivation}} వ్యవస్థ ఉండదు; అంటే, ప్రతి పరిమిత \tetoken{నిర్మాణం}{!!{structure}}~$\Struct{M}$లో
  $\Sat{M}{!B}$ అయితే, అప్పుడే మాత్రమే $\Proves !B$ అయ్యే
  \tetoken{ఒక వ్యుత్పత్తి}{!!a{derivation}} వ్యవస్థ ఉండదు.
\end{cor}

\begin{proof}
  అభ్యాసం.
\end{proof}

\begin{prob}
  \olref[tur][und][tra]{cor:fproof-incomp}ను నిరూపించండి. ప్రతి పరిమిత
  \tetoken{నిర్మాణంలో}{!!{structure}} $!B$ సంతృప్తి చెందితే, అప్పుడే మాత్రమే $\lnot !B$
  పరిమితంగా సంతృప్తిపరచదగినది కాదని గమనించండి. పరిమిత
  సంతృప్తిపరచదగినత, \olref[tur][und][uns]{thm:valid-ce}లోని అర్థంలో ఎందుకు
  అర్ధ-నిర్ణయించదగినదో వివరించండి. పరిమిత చెల్లుబాటుతనానికి
  \tetoken{ఒక వ్యుత్పత్తి}{!!a{derivation}} వ్యవస్థ ఉంటే, పరిమిత సంతృప్తిపరచదగినత నిర్ణయించదగినది
  అయ్యేదని వాదించడానికి దీన్ని ఉపయోగించండి.
\end{prob}

\end{document}
